Fix \(n\ge2\) and the homogeneous policies in \texttt{def:hub-family}.  By \texttt{ass:source-overlap}, \(p(1-p)>0\), so
\[
 \Delta=\frac{q-p}{\sqrt{p(1-p)}}
\]
is well defined.  Since \(p\ne q\), the number \(c=\Delta^2\) is strictly positive and does not depend on \(n\).

For a source Bernoulli coordinate \(Z_j\), define
\[
 \chi_j=\frac{Z_j-p}{\sqrt{p(1-p)}}.
\]
The two possible values of \(Z_j\), together with \texttt{ass:source-product-design}, give
\[
 \mathbb E_{\pi}[\chi_j]
 =\frac{p(1-p)+(1-p)(-p)}{\sqrt{p(1-p)}}=0
\]
and
\[
 \mathbb E_{\pi}[\chi_j^2]
 =\frac{p(1-p)^2+(1-p)p^2}{p(1-p)}=1.
\]
If \(j\ne k\), product-Bernoulli independence and the preceding zero means imply
\[
 \mathbb E_{\pi}[\chi_j\chi_k]
 =\mathbb E_{\pi}[\chi_j]\mathbb E_{\pi}[\chi_k]=0.
\]
These identities also follow directly from the character construction in \texttt{def:local-fourier-character}.

By \texttt{def:hub-family}, the hub support is
\(
 \mathcal S_1=\{\varnothing\}\cup\{\{j\}:j\in V_n\}
\),
whereas for each \(i\ge2\),
\(
 \mathcal S_i=\{\varnothing,\{i\}\}
\).
The spectral-weight formula in \texttt{lem:loomba-eckles-spectral-overlap-envelope}, or equivalently its transcription in \texttt{def:loomba-eckles-spectral-envelope}, therefore gives
\[
 W_1^{\mathrm{LE}}
 =1+\Delta\sum_{j=1}^n\chi_{1,\{j\}}^{\pi},
 \qquad
 W_i^{\mathrm{LE}}
 =1+\Delta\chi_{i,\{i\}}^{\pi}\quad(2\le i\le n).
\]
The moment identities above yield, equation by equation,
\[
\begin{aligned}
 \mathbb E_{\pi}[(W_1^{\mathrm{LE}})^2]
 &=1+2\Delta\sum_{j=1}^n\mathbb E_{\pi}[\chi_j]
   +\Delta^2\sum_{j=1}^n\mathbb E_{\pi}[\chi_j^2]
   +2\Delta^2\sum_{1\le j<k\le n}\mathbb E_{\pi}[\chi_j\chi_k]\\
 &=1+n\Delta^2=1+nc,
\end{aligned}
\]
and, for \(i\ge2\),
\[
 \mathbb E_{\pi}[(W_i^{\mathrm{LE}})^2]
 =1+2\Delta\mathbb E_{\pi}[\chi_i]
   +\Delta^2\mathbb E_{\pi}[\chi_i^2]
 =1+c.
\]

We next verify that these are the exact bounded-class unitwise envelopes, rather than merely upper bounds.  For every \(v_i\in K_i\), \texttt{def:local-evaluation-polytope} gives \(|v_i(x)|\le1\) at every local cell.  Hence, pointwise in the assignment,
\[
 \bigl(v_i(Z_{\Gamma_i})W_i^{\mathrm{LE}}(Z)\bigr)^2
 \le (W_i^{\mathrm{LE}}(Z))^2.
\]
After taking source-design expectation,
\[
 \sup_{v_i\in K_i}
 \mathbb E_{\pi}\!left[
   \bigl(v_i(Z_{\Gamma_i})W_i^{\mathrm{LE}}(Z)\bigr)^2
 \right]
 \le \mathbb E_{\pi}[(W_i^{\mathrm{LE}})^2].
\]
Conversely, the constant schedule \(v_i\equiv1\) is the empty-set character, belongs to every support displayed above, and obeys the pointwise bound.  Thus \(v_i\equiv1\) lies in \(K_i\) and attains equality.  Therefore
\[
 \sup_{v_i\in K_i}
 \mathbb E_{\pi}\!left[
   \bigl(v_i(Z_{\Gamma_i})W_i^{\mathrm{LE}}(Z)\bigr)^2
 \right]
 =\mathbb E_{\pi}[(W_i^{\mathrm{LE}})^2].
\]
Because \(n\ge2\) and \(c>0\), \(1+nc>1+c\).  The definition of the least uniform envelope in \texttt{def:loomba-eckles-spectral-envelope} now gives
\[
 B_n^{\mathrm{LE}}=1+nc.
\]

For the overlap degree, \texttt{def:hub-family} gives \(\Gamma_1=V_n\) and \(\Gamma_i=\{i\}\) for \(i\ge2\).  Thus \(\Gamma_1\cap\Gamma_k=\{k\}\ne\varnothing\) for every \(k\ne1\).  The hub has exactly \(n-1\) intersecting partners, while no unit can have more than \(n-1\) partners.  By \texttt{def:loomba-eckles-spectral-envelope},
\[
 D_n^{\mathrm{LE}}=n-1.
\]
Consequently,
\[
 \frac{B_n^{\mathrm{LE}}(D_n^{\mathrm{LE}}+1)}{n}
 =\frac{(1+nc)n}{n}=1+nc.
\]
Since \(c\in(0,\infty)\) is fixed, this expression is \(\Theta(n)\), and hence also \(\Omega(1)\).  At the smallest admissible instance \(n=2\), the formulas give \(B_2^{\mathrm{LE}}=1+2c\), \(D_2^{\mathrm{LE}}=1\), and the same envelope value \(1+2c\), providing the finite-instance boundary check.  If the excluded boundary \(p=q\) were imposed, then \(c=0\) and the same algebra would give the nonvanishing value one; no division or positivity claim in the proved \(p\ne q\) case uses that boundary.